% Atividade: Análise Combinatória (Fácil)
% Autor: ChatGPT
% Formato: Prova com gabarito comentado no final
\documentclass[12pt,a4paper]{article}
\usepackage[utf8]{inputenc}
\usepackage[brazil]{babel}
\usepackage{amsmath,amssymb}
\usepackage{enumitem}
\usepackage{geometry}
\geometry{margin=2.5cm}
\usepackage{multicol}
\usepackage{titling}
\setlength{\droptitle}{-4em}

\title{Atividade -- Análise Combinatória}
\author{Prof. / Atividade gerada por ChatGPT}
\date{}

\begin{document}
\maketitle

\vspace{0.5cm}
\noindent \textbf{Instruções:} Marque a alternativa correta (A--E). Mostre cálculos quando necessário.\

\begin{enumerate}[leftmargin=*]

\item (Permutação simples) Quantas maneiras diferentes podemos organizar as letras da palavra \textbf{SOL}? 
\begin{multicols}{2}
\begin{enumerate}[label=\Alph*)]
\item 3
\item 4
\item 5
\item 6
\item 8
\end{enumerate}
\end{multicols}

\item (Permutação simples) De quantas formas diferentes podemos organizar 4 pessoas em uma fila?
\begin{multicols}{2}
\begin{enumerate}[label=\Alph*)]
\item 8
\item 10
\item 12
\item 16
\item 24
\end{enumerate}
\end{multicols}

\item (Permutação com repetição) De quantas formas diferentes podemos organizar as letras da palavra \textbf{MAMA}?
\begin{multicols}{2}
\begin{enumerate}[label=\Alph*)]
\item 4
\item 6
\item 8
\item 12
\item 24
\end{enumerate}
\end{multicols}

\item (Arranjo simples) De 5 alunos, queremos escolher 2 para formar uma dupla em ordem (primeiro e segundo lugar). Quantos arranjos possíveis existem?
\begin{multicols}{2}
\begin{enumerate}[label=\Alph*)]
\item 10
\item 15
\item 20
\item 25
\item 30
\end{enumerate}
\end{multicols}

\item (Combinação simples) De 5 alunos, quantos grupos de 2 podemos formar sem importar a ordem?
\begin{multicols}{2}
\begin{enumerate}[label=\Alph*)]
\item 5
\item 8
\item 10
\item 15
\item 20
\end{enumerate}
\end{multicols}

\item (Conceito) Em qual tipo de situação a ordem não importa?
\begin{multicols}{2}
\begin{enumerate}[label=\Alph*)]
\item Permutação
\item Arranjo
\item Combinação
\item Nenhuma
\item Todas
\end{enumerate}
\end{multicols}

\item (Arranjo simples) Em uma corrida com 6 competidores, de quantas maneiras diferentes podemos ter os 3 primeiros colocados?
\begin{multicols}{2}
\begin{enumerate}[label=\Alph*)]
\item 20
\item 60
\item 120
\item 240
\item 720
\end{enumerate}
\end{multicols}

\item (Combinação) Uma professora vai escolher 3 representantes entre 6 alunos. Quantas escolhas diferentes ela pode fazer?
\begin{multicols}{2}
\begin{enumerate}[label=\Alph*)]
\item 10
\item 15
\item 18
\item 20
\item 24
\end{enumerate}
\end{multicols}

\item (Permutação circular) Em uma mesa redonda, 5 pessoas vão se sentar. Quantas disposições diferentes são possíveis?
\begin{multicols}{2}
\begin{enumerate}[label=\Alph*)]
\item 24
\item 60
\item 120
\item 240
\item 720
\end{enumerate}
\end{multicols}

\item (Situação prática – combinação) Uma equipe de esportes precisa escolher 2 jogadores entre 4 para representar o time em uma entrevista. Quantas duplas diferentes podem ser formadas?
\begin{multicols}{2}
\begin{enumerate}[label=\Alph*)]
\item 4
\item 5
\item 6
\item 8
\item 12
\end{enumerate}
\end{multicols}

\end{enumerate}

\newpage
\section*{Gabarito Comentado}
\begin{enumerate}[leftmargin=*]
\item \textbf{D) 6}.\quad Permutação de 3 letras distintas: $3! = 6$.

\item \textbf{E) 24}.\quad Permutação de 4 pessoas em fila: $4! = 24$.

\item \textbf{B) 6}.\quad Permutação com repetição: palavra MAMA tem 4 letras com 2 M e 2 A. Número de arranjos = $\dfrac{4!}{2!\,2!} = \dfrac{24}{4} = 6$.

\item \textbf{C) 20}.\quad Arranjo $A_{5,2} = \dfrac{5!}{(5-2)!} = \dfrac{120}{6} = 20$.

\item \textbf{C) 10}.\quad Combinação $C_{5,2} = \dfrac{5!}{2!\,3!} = 10$.

\item \textbf{C) Combinação}.\quad A combinação é usada quando a ordem dos elementos \emph{não importa}.

\item \textbf{C) 120}.\quad Arranjo $A_{6,3} = \dfrac{6!}{(6-3)!} = \dfrac{720}{6} = 120$.

\item \textbf{D) 20}.\quad Combinação $C_{6,3} = \dfrac{6!}{3!\,3!} = 20$.

\item \textbf{A) 24}.\quad Permutação circular: $(n-1)! = (5-1)! = 4! = 24$.

\item \textbf{C) 6}.\quad Combinação $C_{4,2} = \dfrac{4!}{2!\,2!} = 6$.

\end{enumerate}

\vspace{1cm}
\noindent Boa sorte! Se quiser, posso gerar um PDF pronto para imprimir ou adaptar o estilo (fonte maior, espaço para resolução, incluir cabeçalho para nome/turma etc.).

\end{document}

